\chapter{Relatività speciale}
La fisica nucleare studia le proprietà e le interazioni dei nuclei atomici, dove le energie in gioco possono essere estremamente elevate. In questi contesti, la relatività speciale diventa fondamentale per descrivere correttamente il comportamento delle particelle subatomiche. Ad esempio, la conservazione dell'energia e della quantità di moto deve essere formulata secondo i principi relativistici. Pertanto, lo studio della fisica nucleare è strettamente legato ai concetti fondamentali della relatività speciale.
\section{Trasformazioni di Lorentz}
Le trasformazioni di Lorentz sono un insieme di equazioni che descrivono come le coordinate spazio-temporali di un evento cambiano quando si passa da un sistema di riferimento inerziale a un altro che si muove a velocità costante rispetto al primo.\\
Scrivendo le trasformazioni di Lorentz per un sistema che si muove lungo l'asse $z$ con velocità $v$, abbiamo:
\begin{equation}
    \begin{cases}
        t' = \gamma \left( t - \dfrac{v z}{c^2} \right) \\
        x' = x \\
        y' = y \\
        z' = \gamma (z - v t)
    \end{cases}
\end{equation}
dove $\gamma = \dfrac{1}{\sqrt{1 - \dfrac{v^2}{c^2}}}$ è il fattore di Lorentz e $c$ è la velocità della luce nel vuoto.\\
Formulando le trasformazioni di Lorentz lungo l'asse $z$, per le trasformazioni delle velocità otteniamo:
\begin{equation}
    \begin{cases}
        u'_x = \dfrac{u_x}{\gamma \left( 1 - \dfrac{v u_z}{c^2} \right)} \\
        u'_y = \dfrac{u_y}{\gamma \left( 1 - \dfrac{v u_z}{c^2} \right)} \\
        u'_z = \dfrac{u_z - v}{1 - \dfrac{v u_z}{c^2}} \\
    \end{cases}
\end{equation}
Infine, la matrice di Lorentz $\Lambda$ risulta essere:
\begin{equation}
    \Lambda =
    \begin{pmatrix}
        \gamma & 0 & 0 & -\beta \gamma  \\
        0 & 1 & 0 & 0 \\
        0 & 0 & 1& 0\\
        -\beta \gamma & 0 & 0 & \gamma \\
    \end{pmatrix}
\end{equation}
dove $\beta = \dfrac{v}{c}$.\\
Dalle trasformazioni di Lorentz otteniamo risultati non previsti dalla fisica classica come la \textbf{contrazione delle lunghezze} e la \textbf{dilatazione dei tempi}:
\begin{itemize}
    \item $L(v) = \dfrac{L'}{\gamma(v)}$
    \item $\Delta t (v) = \Delta t' \gamma (v)$
\end{itemize}
Le quantità $L'$ ed $\Delta t'$ sono dette rispettivamente lunghezza propria e intervallo di tempo proprio e sono misurate nel sistema di riferimento inerziale nel quale l'oggetto è stazionario.
\section{Evidenze scientifiche}
\subsection{I muoni}
Grazie alla teoria della relatività speciale riusciamo a descrivere diversi fenomeni incompatibili con la meccanica classica.
Prendiamo ad esempio la vita dei raggi cosmici, radiazioni formate all'$86\%$ da protoni, al $13\%$ da elio, al $2\%$ da elettroni, all'$1\%$ da altri nuclei e da qualche raggio $\gamma$.\\
Questi raggi sono originati da diverse entità presenti nello spazio come stelle di varia tipologia, supernove o stelle di neutroni, anche dette stelle compatte a causa della loro elevata densità.\\
Secondo la meccanica classica alcune particelle originate dall'impatto dei raggi cosmici con la nostra atmosfera, i \textbf{muoni} $\mu$, non dovrebbero essere in grado di arrivare sulla superficie terrestre a causa del loro breve tempo di vita.

\addfigure[Raggi cosmici in atmosfera. Fonte: \href{https://www.vox.com/the-highlight/2019/7/16/17690740/cosmic-rays-universe-theory-science}{Vox, Nasa}]{png/cosmic_rays_diagram}{0.6}

Tuttavia, grazie alla relatività speciale riusciamo a calcolare la dilatazione dei tempi che interessa queste particelle, in particolare il fattore $\gamma(v)$ è pari a circa 30.
Se ci limitiamo alla meccanica classica, con un tempo di vita medio di $\tau = \qty{1.56}{\micro\second}$ il rapporto di sopravvivenza, che si ottiene confrontando le misurazioni effettuate a $\qty{10}{\kilo\metre}$ di altezza con quelle effettuate sulla superficie, dovrebbe essere:
\begin{equation}
    \dfrac{I}{I_0} \simeq 0.26 \cdot 10^{-6}
\end{equation}
Al contrario, sperimentalmente, osserviamo che il rapporto di sopravvivenza è:
\begin{equation}
    \dfrac{I}{I_0} \simeq 0.049
\end{equation}
Dato compatibile con le previsioni della relatività speciale.

\subsection{Cutoff GZK}

\addfigure[Energie e rilevamento dei raggi cosmici. Fonte: \href{https://www.researchgate.net/publication/360295707_A_Search_for_the_Origin_of_Ultra-High_Energy_Neutrinos_with_ANITA-4}{ResearchGate}]{jpg/GZK_cutoff}{0.7}

Nell'osservazione dei raggi cosmici possiamo individuare un andamento delle possibili energie assunte in particolar modo dai protoni. Come prevedibile all'aumentare dell'energia diminuisce il numero di protoni rilevati, tuttavia osserviamo una discesa netta precedente ad una certa soglia energetica $\simeq \qty{e20}{\electronvolt}$. Per spiegare questo fenomeno occorre avere in mente il modello cosmologico, che possiamo riassumere con i seguenti punti fondamentali:
\begin{itemize}
  \item Presenza di un Big Bang $\simeq 13.8 \cdot 10^9$ anni fa
  \item Espansione dell'universo e abbassamento della temperatura
\end{itemize}
Questa teoria è giustificata da un'importante osservazione sperimentale, ovvero quella della \textbf{radiazione cosmica di fondo} (CMB). Questa radiazione permea lo spazio ed assume uno spettro essenzialmente identico a quello di un corpo nero alla temperatura di $\qty{2.725 \pm 0.001}{\kelvin}$.

\addsvg[Spettro della radiazione cosmica di fondo. Fonte: \href{https://en.wikipedia.org/wiki/Cosmic_microwave_background}{Wikipedia}]{Cmbr}{0.8}

Tornando dunque al problema della rilevazione dei raggi cosmici, possiamo ipotizzare che dei protoni ad alta energia siano in grado di interagire con i fotoni della radiazione cosmica di fondo perdendo energia e generando un pione:
\begin{equation}
  p + \gamma_{CMB} \rightarrow p + \pi^0
\end{equation}
Per calcolare l'energia minima alla realizzazione di questo processo poniamo che il protone e il pione siano entrambi a riposo dopo l'interazione e che l'impatto avvenga lungo un'unica direzione ovvero:
\begin{equation}
  \begin{split}
    &p_{p}^\mu = \brackets{E_p, 0, 0, E_p} \\[8pt]
    &p_{\pi}^\mu = \brackets{E_{\text{CMB}}, 0, 0, -E_{\text{CMB}}} \\[8pt]
    &p_{p'}^\mu = \brackets{m_{p}, 0, 0, 0} \\[8pt]
    &p_{\pi}^\mu = \brackets{m_\pi, 0, 0, 0}
  \end{split}
\end{equation}
Dove abbiamo utilizzato l'approssimazione ultra-relativistica $E \simeq K \simeq p$. Dalla conservazione del modulo del tensore otteniamo che:
\begin{equation}
  E_p = \dfrac{2m_p m_{\pi} + m^2_\pi}{4E_{\text{CMB}}}
\end{equation}
Dai dati sperimentali sappiamo che:
\begin{itemize}
  \item $E_{\text{CMB}} \simeq \qty{3e-4}{\electronvolt}$
  \item $m_p = \qty{938}{\mega\electronvolt}$
  \item $m_\pi = \qty{135}{\mega\electronvolt}$
\end{itemize}
In questo modo troviamo:
\begin{equation}
  E_p \simeq \qty{2e20}{\electronvolt} \simeq \qty{30}{\joule}
\end{equation}
Questo spiega l'andamento qualitativo delle energie dei protoni, tuttavia, tenendo in considerazione la natura termica della CMB la soglia GKZ viene stimata a circa $\qty{5e19}{\electronvolt}$, dunque ci sono dei raggi cosmici che violano questo limite. La spiegazione completa di questo fenomeno è ancora oggi oggetto di studio.
\section{Formalismo tensoriale}
Il formalismo tensoriale è uno dei principali strumenti della relatività speciale, in quanto permette di descrivere le grandezze fisiche mediante oggetti matematici chiamati tensori. I tensori sono generalizzazioni dei vettori e delle matrici. Un tensore di ordine 1 ad esempio è indicato da un solo indice superiore, se contravarianti, o inferiore, se covarianti. Un tensore di ordine 2 invece è indicato da due indici, che possono essere entrambi superiori, entrambi inferiori o uno superiore e uno inferiore.\\
Un esempio di tensore di ordine 1 è il quadrivettore posizione $x^\mu = (ct, x, y, z)$, mentre un esempio di tensore di ordine. Utilizzeremo la notazione di Einstein, che prevede la somma implicita sugli indici ripetuti, per esempio:
\begin{equation}
    A^\mu B_\mu = \sum_{\mu=0}^{3} A^\mu B_\mu
\end{equation}
Dunque le trasformazioni di Lorentz in notazione tensoriale si scrivono come:
\begin{equation}
    x'^\mu = \Lambda^\mu_{\ \nu} x^\nu
\end{equation}
Un altro concetto importante nel formalismo tensoriale è il tensore metrico $g_{\mu \nu}$, che permette di alzare e abbassare gli indici dei tensori. In relatività speciale, il tensore metrico è dato da:
\begin{equation}
    g_{\mu \nu} =
    \begin{pmatrix}
        1 & 0 & 0 & 0 \\
        0 & -1 & 0 & 0 \\
        0 & 0 & -1 & 0 \\
        0 & 0 & 0 & -1 \\
    \end{pmatrix}
\end{equation}
In questo caso abbiamo adottato la convenzione del segno $(+,-,-,-)$. Utilizzando il tensore metrico, possiamo scrivere la relazione tra un tensore covariante e uno contravariante come:
\begin{equation}
    A_\mu = g_{\mu \nu} A^\nu
\end{equation}
Poiché il formalismo tensoriale e la definizione dei quadrivettori discende dalle trasformate di Lorentz e dunque dall'invarianza della velocità della luce abbiamo che la quantità:
\begin{equation}
    x^\mu x_\mu = s^2 = c^2 t^2 - x^2 - y^2 - z^2
\end{equation}
È invariante sotto trasformazioni di Lorentz. Inoltre qualsiasi prodotto scalare tra due quadrivettori è invariante, ad esempio:
\begin{equation}
    p^\mu p_\mu = \dfrac{E^2}{c^2} - \vec{p}^2 = m^2 c^2 \implies E^2 = \vec{p}^2c^2 + m^2 c^4
\end{equation}
In questo modo otteniamo la formula fondamentale che lega energia, quantità di moto e massa a riposo di una particella. Utilizzeremo la notazione dove $m$ è la massa a riposo, mentre $m(v)$ è la massa relativistica, che dipende dalla velocità della particella:
\begin{equation}
    m(v) = \gamma(v) m
\end{equation}
\section{Dinamica relativistica}
Per descrivere i processi dinamici relativistici occorre considerare la conservazione di:
\begin{itemize}
    \item Energia $E$
    \item Quantità di moto $\vec{p}$
\end{itemize}
Dunque possiamo affermare che tutte le componenti del tensore $p^\mu$ sono conservate. Inoltre sappiamo che l'energia può essere espressa come:
\begin{equation}
    E = K + U
\end{equation}
Sapendo che $U = mc^2$ è l'energia a riposo abbiamo che:
\begin{equation}
    \begin{split}
        &E^2 = \vec{p}^2c^2 + m^2c^4 \\[8pt]
        &\brackets{K + mc^2}^2 = \vec{p}^2c^2 + m^2c^4 \\[8pt]
        &K^2 + 2mKc^2 + \cancel{m^2c^4} = \vec{p}^2c^2 + \cancel{m^2c^4} \\[8pt]
        &K^2 + 2mKc^2 - \vec{p}^2c^2 = 0
    \end{split}
\end{equation}
Con questa espressione possiamo individuare tre diverse regioni in base all'energia cinetica della particella che stiamo considerando:
\begin{itemize}
    \item Regione non-relativistica:
    \begin{equation}
        \begin{split}
            &K \ll mc^2 \implies K^2 + 2mKc^2 - \vec{p}^2c^2 \simeq 2mKc^2 - \vec{p}^2c^2 \\[8pt]
            &2mKc^2 - \vec{p}^2c^2 \simeq 0 \\[8pt]
            &K \simeq \dfrac{\vec{p}^2}{2m}
        \end{split}
    \end{equation}
    \item Regione relativistica:
    \begin{equation}
        K \sim mc^2
    \end{equation}
    \item Regione ultra-relativistica (U.R.):
    \begin{equation}
        \begin{split}
            &K \gg mc^2 \implies K^2 + 2mKc^2 - \vec{p}^2c^2 \simeq K^2 - \vec{p}^2c^2 \\[8pt]
            &K^2 - \vec{p}^2c^2 \simeq 0 \\[8pt]
            &K^2 \simeq \vec{p}^2c^2
        \end{split}
    \end{equation}
\end{itemize}
Per la regione ultra-relativistica possiamo notare che l'approssimazione diventa esatta nel caso di particelle con massa nulla, ovvero i fotoni. Infatti possiamo scrivere l'energia di un fotone come:
\begin{equation}
    E_\gamma = pc
\end{equation}
Possiamo ricavare questa formula dall'equazione generale dell'energia utilizzata in relatività speciale:
\begin{equation}
    E = \pm \sqrt{m^2c^4 + \vec{p}^2c^2}
\end{equation}
Ponendo $m = 0$ otteniamo la formula precedentemente citata. Inoltre osserviamo che i valori dell'energia ottenuti dalla formula possono essere sia positivi che negativi. Potremmo essere tentati di scartare a priori i valori negativi etichettandoli come ``non fisici'', tuttavia torneremo sul significato di energie negative successivamente.\\
\textbf{Esempio}\\
Ipotizziamo di avere una particella a riposo $\pi^0$ con massa $m_\pi c^2= \qty{135}{\mega\electronvolt}$. La particella decade in due raggi $\gamma$:
\begin{equation}
  \pi^0 \rightarrow \gamma + \gamma
\end{equation}
Ponendoci nel sistema di riferimento stazionario per $\pi^0$ abbiamo che:
\begin{equation}
  p_{\pi}^\mu = \brackets{m_\pi, 0, 0, 0}
\end{equation}
Poiché il tensore $p^\mu$ deve essere conservato abbiamo:
\begin{equation}
  \begin{split}
    &p_{\gamma_1}^\mu = \brackets{E_{\gamma_1}, 0, 0, p_{z,\gamma_1}} \\[8pt]
    &p_{\gamma_2}^\mu = \brackets{E_{\gamma_2}, 0, 0, p_{z,\gamma_2}}
  \end{split}
\end{equation}
Dove abbiamo scelto l'asse $z$ parallelo alla direzione dei fotoni. Abbiamo dunque che:
\begin{equation}
  \begin{split}
    E_{\gamma_1} + E_{\gamma_2} = m_\pi
    p_{z,\gamma_1} = - p_{z,\gamma_2}
  \end{split}
\end{equation}
Dunque abbiamo:
\begin{equation}
  \begin{split}
    &p_{\gamma_1}^\mu = \brackets{E, 0, 0, p_z} \\[8pt]
    &p_{\gamma_2}^\mu = \brackets{m_\pi - E, 0, 0, -p_z}
  \end{split}
\end{equation}
Sapendo che $E_\gamma = p$ abbiamo che:
\begin{equation}
  2p = m_\pi \implies p = \dfrac{m_\pi}{2}
\end{equation}
Dunque:
\begin{equation}
  \begin{split}
    &p_{\gamma_1}^\mu = \brackets{\dfrac{m_\pi}{2}, 0, 0, \dfrac{m_\pi}{2}} \\[8pt]
    &p_{\gamma_2}^\mu = \brackets{m_\pi - \dfrac{m_\pi}{2}, 0, 0, -\dfrac{m_\pi}{2}}
  \end{split}
\end{equation}
Dunque:
\begin{equation}
  \begin{split}
    &p_{\gamma_1}^\mu = \brackets{\dfrac{m_\pi}{2}, 0, 0, \dfrac{m_\pi}{2}} \\[8pt]
    &p_{\gamma_2}^\mu = \brackets{\dfrac{m_\pi}{2}, 0, 0, -\dfrac{m_\pi}{2}}
  \end{split}
\end{equation}
